% !TEX root = ../main.tex
\chapter{Institutional Warrant and Tiered Reliance}
\label{ch:warrant-tiers}
\label{sec:tiers}

The architecture of Section~\ref{sec:governed-rag} serves a condition I have so far
used without defining. This section defines it and makes the controls proportional
to consequence.

\section{Warrant}

Grounding is necessary for reliance and not sufficient. An assertion can rest on
authentic, applicable, authoritative evidence, be traceable to it, and still be
something the institution is not entitled to act on---because nobody checked it, or
because nobody with the standing to authorize reliance did so.

\begin{definition}[Institutional warrant]
\label{def:warrant}
\begin{equation}
  \begin{aligned}
  \Warranted_I(p \mid c) \;\leftrightarrow\;
    &\Resolved(c) \wedge \Grd_S(p \mid c)\\
    &\wedge\ \Verified(V, p, c) \wedge \Authorized(I, p, c).
  \end{aligned}
  \label{eq:warrant}
\end{equation}
\end{definition}

where

\begin{itemize}
  \item $\Resolved(c)$: the material context variables of \eqref{eq:context} are
        explicit rather than assumed or inferred without validation;
  \item $\Grd_S(p \mid c)$: Definition~\ref{def:grounding}---which unpacks to the
        six evidence conditions of \eqref{eq:grounded} plus $\Sensitive$, so
        authenticity, versioning, authority, and applicability are preconditions of
        valid support rather than separate requirements;
  \item $\Verified(V, p, c)$: an appropriate procedure has checked the assertion,
        its evidence relation, and the controls the decision class requires;
  \item $\Authorized(I, p, c)$: the institution permits reliance at the relevant
        risk tier, possibly subject to human review or approval.
\end{itemize}

Three properties of \eqref{eq:warrant} matter for how it is used.

It is a conjunction, not a score. There is no aggregation of the four conjuncts into
a confidence figure, and an assertion satisfying three of them is not
three-quarters warranted. This is the same point as the flow in
Section~\ref{sec:governed-rag}: these are distinct conditions with distinct owners
and distinct failure records.

It is a property of a process, not of a model. Nothing in \eqref{eq:warrant} is a
property of model weights, a benchmark score, or a generated response.
$\Authorized$ in particular is a fact about an institution's governing rules and
delegations, and no improvement to $M$ can supply it. This is why warrant cannot be
procured: a vendor can sell you a component that makes grounding measurable, and
cannot sell you an authorization.

It is revocable and context-relative. \eqref{eq:warrant} is not a claim that $p$ is
certain, or even that $p$ is true. It is the condition under which an institution
may responsibly act on $p$ for a defined purpose, knowing that it may be wrong and
that the record will show what was relied upon and why. That is the standard
ordinary institutional decision-making already meets, and the standard machine
assertion currently does not.

\section{Reliance tiers}

Controls should be proportional to the consequence of reliance. The table below is
deliberately compact: its purpose is to show that permissible reliance is set by
consequence rather than by model quality, and that argument does not require a
fully specified control catalog. Any real program will need its own tiering,
calibrated to its regulator and its risk appetite.

\begin{table}[ht]
  \centering
  \small
  \begin{tabular}{@{}>{\raggedright\arraybackslash}p{1.1cm}>{\raggedright\arraybackslash}p{3.6cm}>{\raggedright\arraybackslash}p{4.4cm}>{\raggedright\arraybackslash}p{4.4cm}@{}}
    \toprule
    Tier & Use & Required & Invariants engaged \\
    \midrule
    0 & Unwarranted generative assistance: brainstorming, drafting
      & Clear framing as unverified; no reliance claim
      & none \\
    \addlinespace
    1 & Informational assistance
      & Visible uncertainty; no reliance claim
      & I17 \\
    \addlinespace
    2 & Evidence-linked assistance
      & Sources visible; user retains verification duty
      & I8, I9, I12 \\
    \addlinespace
    3 & Governed decision support
      & Context-bound evidence, verification controls, audit trail, authorized
        human reliance
      & I8--I19 \\
    \addlinespace
    4 & Constrained execution
      & Deterministic guardrails, approvals, reversibility, monitoring, full
        records
      & I8--I19 plus action controls \\
    \addlinespace
    5 & Autonomous consequential action
      & Narrowly specified settings, stringent formal controls, continuing
        assurance, explicit authority to act
      & all, plus independent assurance \\
    \bottomrule
  \end{tabular}
  \caption{Reliance tiers. The tier is set by the consequence of reliance, not by
    model quality. Tier~5 is included for completeness and should be approached
    with considerable caution in regulated settings; nothing in this book argues
    that current systems are ready for it.}
  \label{tab:tiers}
\end{table}

Assignment to a tier needs criteria, and while a full consequence taxonomy is out of
scope, the dimensions are nameable: \emph{reversibility} of the action the assertion
supports; \emph{magnitude} of the exposure if it is wrong; \emph{breadth}, meaning how
many parties are affected and whether they consented; \emph{detectability}, whether an
error would surface before harm; and \emph{availability of a competent human check} in
the time the decision allows. A use case scoring badly on reversibility and
detectability belongs at a higher tier than its nominal risk rating suggests, because
those two dimensions determine whether anything can be done after the fact.

The principle the table exists to enforce is a negative one. Permissible reliance
does not rise because the model got better. It rises when evidence, verification,
constraints, authorization, and accountability rise together and in proportion to
what is at stake. A system that improves its benchmark score has not thereby earned
a higher tier, and a system that has implemented I8--I19 has not thereby earned
Tier~5.

Two boundaries are worth marking explicitly, because they are where programs
drift. The step from Tier~2 to Tier~3 is the step at which the institution rather
than the user assumes the verification duty, and it requires everything in
Table~\ref{tab:invariants-system}; systems are frequently described as Tier~3 while
implementing Tier~2. And the step from Tier~3 to Tier~4 is the step from assertion
to action, at which reversibility and blast radius become first-order design
questions rather than operational details.

\section{Three Teams, Three Budgets}
\label{sec:three-teams}

Everything above is stated as though an organization were a single agent that could
decide to satisfy Table~\ref{tab:invariants-system}. It is not, and the way it is
divided is itself a principal cause of the failures this book describes. This
section is the least formal in the book and, for practitioners, possibly the most
useful.

\subsection{The division}

Responsibility for the conditions of warrant is distributed across organizational
functions that were not designed to hold it jointly. In financial services the fourth
row below is a second-line function with an independent reporting line and a
supervisory mandate; in other sectors the equivalent may sit in quality assurance,
clinical governance, or internal audit. Wherever it sits, it is the function that
performs the independent challenge this book argues for, and it is the one most often
absent from discussions of these systems.

\begin{table}[ht]
  \centering
  \small
  \begin{tabular}{@{}>{\raggedright\arraybackslash}p{2.8cm}>{\raggedright\arraybackslash}p{4.6cm}>{\raggedright\arraybackslash}p{2.4cm}>{\raggedright\arraybackslash}p{4.0cm}@{}}
    \toprule
    Function & Typical responsibility & Owns & Failure if left in isolation \\
    \midrule
    Data / knowledge
      & Source ingestion, curation, lineage, metadata, quality
      & I8, I9
      & Authentic documents with no policy applicability or output control \\
    \addlinespace
    Policy / risk / compliance
      & Rules, jurisdiction, decision authority, approvals, procedures
      & I10, I15
      & Governing policy with no control over how assertions are retrieved,
        generated, or verified \\
    \addlinespace
    Platform / AI engineering
      & Models, prompts, orchestration, retrieval, evaluation, monitoring
      & I11--I14, I17
      & Answer quality, latency, and usability with no evidentiary authority or
        institutional authorization \\
    \addlinespace
    Model risk / independent validation
      & Independent challenge, validation, and approval of models and their use
      & I13, I15, I16
      & Validating $M$ while the risk sits in $S$: a well-documented model, an
        unexamined evidence path \\
    \bottomrule
  \end{tabular}
  \caption{Four functions, and the invariants each is positioned to supply. No
    function owns a complete path from source to authorized reliance.}
  \label{tab:three-teams}
\end{table}

These are not merely three technical workstreams. They typically have separate
budgets, separate incentives, separate release cadences, separate accountability
structures, and separate systems of record. A data team is measured on lineage
coverage and freshness. A compliance function is measured on policy completeness
and audit findings. A platform team is measured on answer quality, latency,
adoption, and cost. None is measured on the grounding rate of
\eqref{eq:rate}, because it spans all three.

\subsection{The unit of failure is the handoff}

The consequence is that failures concentrate at the boundaries rather than inside
any function's remit. A system may have:

\begin{itemize}
  \item excellent data lineage and weak policy binding---I8 and I9 without I10;
  \item well-specified governing policy and uncontrolled retrieval---I10 without
        I11;
  \item high benchmark accuracy and no per-assertion provenance---nothing in
        Table~\ref{tab:invariants-system} at all, with a validation report that
        looks complete;
  \item citations without applicability validation---I12 without I11, which is
        precisely cell (v);
  \item human review without a reproducible evidence path---review that cannot be
        re-performed and therefore cannot be audited, which is I13 without I14;
  \item a validated model inside an unvalidated system---documentation of $M$'s
        training, evaluation, and limitations, with no examination of context
        resolution, applicability logic, defeater checking, or authorization. This is
        the characteristic failure of the fourth row, and it is the most dangerous
        because it produces a formal approval.
\end{itemize}

Each of these is a system in which every function has done its job. The ungoverned
handoff is the operational unit of failure, and it is invisible from inside any one
function's metrics.

This also explains a pattern worth naming, because it is the most expensive form of
the problem. Cell (v)---misapplied grounding---sits exactly on the boundary between
the data function, which supplies authentic versioned sources, and the policy
function, which knows which sources govern which decisions. The metadata required to
enforce applicability is a policy artifact stored in a data system and consumed by a
platform component. It is nobody's deliverable, and it is the single most common
thing I find missing.

\subsection{What follows}

Two conclusions, and the first is uncomfortable for how these programs are
usually organized.

Evidentiary controls must be architecture and cross-functional invariants, not
after-the-fact review tasks. A review step owned by one function cannot establish a
property that requires artifacts from two others, and adding reviewers does not
change this. Table~\ref{tab:invariants-system} is a specification for a system, and
a specification that spans three budgets needs an owner who spans them.

A fifth function appears once retention is considered. Privacy and records
management own the tension described in Section~\ref{sec:retention}: the audit record
that reconstructability requires is in direct conflict with minimization and erasure
obligations, and neither the data, policy, platform, nor validation function is
positioned to resolve it. That handoff has the same structure as the others and is
usually discovered later.

And the grounding rate is the right cross-functional metric precisely because no
single function can move it alone. Accuracy can be improved by the platform team
without reference to the others, which is part of why it is the number that gets
reported. \eqref{eq:rate} cannot: improving it requires corpus work, policy
metadata, and pipeline control together. A metric that only improves when three
functions cooperate is an unusually honest instrument for a problem whose failures
live between them.
